% Chapter 02.8: নেটওয়ার্ক ডিভাইস
\begingroup
\chTwoReadableTopicStyle
\section{নেটওয়ার্ক ডিভাইস}

\subsection{বাস্তব সমস্যা ও প্রয়োজন}

একটি কম্পিউটার ল্যাবের সব কম্পিউটার নেটওয়ার্কে যুক্ত করতে \textbf{সুইচ} ব্যবহার করা হয়। ল্যাবের নেটওয়ার্ককে ইন্টারনেটে যুক্ত করতে \textbf{রাউটার} ও \textbf{মডেম} ব্যবহার করা হতে পারে। তারবিহীন ইন্টারনেট সেবা দিতে \textbf{অ্যাকসেস পয়েন্ট} বা Wi-Fi রাউটার প্রয়োজন হয়।

কম্পিউটার, সার্ভার, প্রিন্টার, ইন্টারনেট সংযোগ ও তারবিহীন ডিভাইস যুক্ত করতে বিভিন্ন হার্ডওয়্যার যন্ত্র প্রয়োজন। এসব যন্ত্র সংযোগ নিয়ন্ত্রণ, সংকেতের শক্তি বাড়ানো, ডেটা ফরওয়ার্ড করা ও নিরাপত্তা নিয়ন্ত্রণেও সাহায্য করে। এগুলোকে \textbf{নেটওয়ার্ক ডিভাইস} বলে।

\begin{keypointbox}{মূল ধারণা}
নেটওয়ার্ক ডিভাইসের কাজ বুঝতে চারটি বিষয় দেখতে হয়:
\begin{itemize}
\item ডিভাইসটি নেটওয়ার্কের ভেতরে কাজ করে, নাকি ভিন্ন নেটওয়ার্ক সংযুক্ত করে?
\item ডিভাইসটি সংকেতের শক্তি বাড়ায়, নাকি ডেটা সঠিক গন্তব্যে পাঠায়?
\item ডিভাইসটি তারযুক্ত সংযোগ দেয়, নাকি বেতার সংযোগ দেয়?
\item ডিভাইসটি নিরাপত্তা নিয়ন্ত্রণ বা ট্রাফিক বাছাই করে কি না?
\end{itemize}
\end{keypointbox}

\subsection{নেটওয়ার্ক ডিভাইস}

\begin{definitionbox}{নেটওয়ার্ক ডিভাইস}
কম্পিউটার, সার্ভার, প্রিন্টার ও বিভিন্ন নেটওয়ার্ক সংযুক্ত করা, ডেটা আদান-প্রদান নিয়ন্ত্রণ, সংকেত পুনর্গঠন, পথ নির্ধারণ বা নিরাপত্তা নিয়ন্ত্রণে ব্যবহৃত হার্ডওয়্যার কিংবা হার্ডওয়্যার-সফটওয়্যার সমন্বিত উপাদানকে নেটওয়ার্ক ডিভাইস বলে।
\end{definitionbox}

নেটওয়ার্ক ডিভাইসকে Networking Hardware-ও বলা হয়। LAN, MAN, WAN, ইন্টারনেট সংযোগ বা তারবিহীন নেটওয়ার্ক কার্যকরভাবে পরিচালনায় এসব যন্ত্র প্রয়োজন।

\begin{examplebox}{সহজ উদাহরণ}
\begin{itemize}
\item কম্পিউটার ল্যাবে অনেক PC যুক্ত করতে সুইচ ব্যবহার করা হয়।
\item বাসার Wi-Fi নেটওয়ার্কে রাউটার একাধিক যন্ত্রকে ইন্টারনেট ব্যবহারের সুযোগ দেয়।
\item ISP-এর উপযুক্ত সংযোগ থেকে ডিজিটাল ডেটা নিতে মডেম ব্যবহার করা হয়।
\item দীর্ঘ কেবলপথে দুর্বল সংকেত পুনর্গঠন করে আবার পাঠাতে রিপিটার ব্যবহার করা হয়।
\item ভিন্ন প্রোটোকল বা নেটওয়ার্ক ব্যবস্থা যুক্ত করতে গেটওয়ে ব্যবহার করা হয়।
\end{itemize}
\end{examplebox}


\chTwoImage{chapters/Chapter02/images/network_devices_overview.png}{নেটওয়ার্ক ডিভাইসের সারসংক্ষেপ}


\subsection{Network Device কেন প্রয়োজন?}

নেটওয়ার্কে ডেটা কেবল দিয়ে চলাচল করলেই সব কাজ শেষ হয় না। কোন ডিভাইস বা নেটওয়ার্কে ডেটা পাঠাতে হবে, দুর্বল সংকেত কীভাবে পুনর্গঠন করতে হবে এবং অননুমোদিত প্রবেশ কীভাবে নিয়ন্ত্রণ করতে হবে—এসব কাজে নেটওয়ার্ক ডিভাইস প্রয়োজন।

\begin{keypointbox}{Network device-এর প্রয়োজনীয়তা}
\begin{itemize}
\item একাধিক কম্পিউটার বা ডিভাইস নেটওয়ার্কে যুক্ত করে।
\item ডেটা প্যাকেট সঠিক গন্তব্যে পৌঁছাতে সাহায্য করে।
\item স্থানীয় নেটওয়ার্ককে ইন্টারনেট বা অন্য নেটওয়ার্কের সঙ্গে যুক্ত করে।
\item যন্ত্রের কাজ অনুযায়ী দুর্বল সংকেত পুনর্গঠন বা বর্ধন করতে পারে।
\item বেতার ডিভাইসকে নেটওয়ার্কে যুক্ত হওয়ার সুযোগ দেয়।
\item নেটওয়ার্ক ট্রাফিক নিয়ন্ত্রণ করে কার্যক্ষমতা বাড়াতে সাহায্য করে।
\item অননুমোদিত প্রবেশ ও ক্ষতিকর ট্রাফিক বাছাই করে বাধা দিতে পারে।
\end{itemize}
\end{keypointbox}

\subsection{OSI Layer অনুযায়ী Network Device}

\chTwoPair{ডিভাইস কোন তথ্য ব্যবহার করে সিদ্ধান্ত নেয়, তা OSI স্তরের সঙ্গে মিলিয়ে বোঝা যায়। NIC-এ ভৌত ইন্টারফেস ও ডেটা লিংক স্তরের কাজ দুটিই আছে। Multilayer Switch একাধিক স্তরে কাজ করতে পারে। গেটওয়ের স্তর নির্ভর করে রূপান্তরের ধরনের ওপর।

}{\chTwoImage{chapters/Chapter02/images/osi_layer_network_devices_map.png}{OSI স্তর অনুযায়ী নেটওয়ার্ক ডিভাইস}}
\begin{center}
\begin{tabular}{|p{\dimexpr 0.2714\linewidth-2\tabcolsep-\arrayrulewidth\relax}|p{\dimexpr 0.3844\linewidth-2\tabcolsep-\arrayrulewidth\relax}|p{\dimexpr 0.3392\linewidth-2\tabcolsep-\arrayrulewidth\relax}|}
\hline
\textbf{OSI Layer} & \textbf{Device} & \textbf{প্রধান কাজ} \\
\hline
Physical Layer & Repeater, Hub, Modem & সংকেত পাঠানো, পুনর্গঠন বা রূপান্তর করা \\
\hline
Data Link Layer & NIC, Bridge, Switch & ফ্রেম তৈরি ও গ্রহণ; ব্রিজ ও সুইচে MAC-ভিত্তিক ফরওয়ার্ডিং \\
\hline
Network Layer & Router & IP Address অনুযায়ী প্যাকেটের পথ নির্ধারণ \\
\hline
Multiple Layer & Gateway, Firewall & প্রোটোকল রূপান্তর, ট্রাফিক বাছাই ও নিরাপত্তা \\
\hline
\end{tabular}
\end{center}





\subsection{Network Interface Card বা NIC}

\begin{definitionbox}{NIC}
কম্পিউটার বা অন্য ডিভাইসকে নেটওয়ার্কে সংযুক্ত করার হার্ডওয়্যার উপাদানকে Network Interface Card বা NIC বলে।
\end{definitionbox}

প্রচলিত Ethernet MAC (Media Access Control) Address ৪৮ bit বা ৬ Byte দীর্ঘ। এটি ১২টি হেক্সাডেসিমাল অঙ্ক দিয়ে লেখা যায়; যেমন 02:1A:2B:3C:4D:5E। MAC ডেটা লিংক স্তরের ঠিকানা এবং IP Address নেটওয়ার্ক স্তরের ঠিকানা।
\textbf{IP Address-এর দৈর্ঘ্য:} IPv4 ঠিকানা ৩২ bit বা ৪ Byte এবং IPv6 ঠিকানা ১২৮ bit বা ১৬ Byte দীর্ঘ। MAC ও IP ঠিকানার বিটসংখ্যা আলাদা। রাউটার কোন ধরনের ঠিকানা ব্যবহার করে, তার সঙ্গে দৈর্ঘ্যের তথ্য মিলিয়ে বুঝতে হয়।
\textbf{ফ্রেমের ঠিকানার উৎস:} হোস্টের নেটওয়ার্ক সফটওয়্যার ও NIC ফ্রেম প্রস্তুত করে। Source MAC ওই লিংকের প্রেরক ইন্টারফেসের ঠিকানা। Destination MAC একই লিংকের প্রাপক বা পরবর্তী ধাপের ইন্টারফেসের ঠিকানা। সুইচ এই ঠিকানা ব্যবহার করে। ঠিকানা শেখার ধাপ সুইচের অংশে আছে।


NIC কম্পিউটারকে নেটওয়ার্কের মাধ্যমের সঙ্গে যুক্ত করে। প্রতিটি NIC-এর সাধারণত একটি স্বতন্ত্র ভৌত ঠিকানা থাকে, যার নাম MAC Address। তারযুক্ত NIC Ethernet কেবল এবং তারবিহীন NIC Wi-Fi সংকেত ব্যবহার করে।

\chTwoPair{\begin{keypointbox}{NIC-এর বৈশিষ্ট্য}
\begin{itemize}
\item কম্পিউটার বা অন্য ডিভাইসকে নেটওয়ার্কে সংযুক্ত করে।
\item তারযুক্ত বা তারবিহীন—উভয় ধরনের হতে পারে।
\item MAC Address দিয়ে সংশ্লিষ্ট নেটওয়ার্ক ইন্টারফেস শনাক্ত করা যায়।
\item LAN-এর যোগাযোগে গুরুত্বপূর্ণ ভূমিকা রাখে।
\item মাদারবোর্ডে যুক্ত থাকতে পারে বা পৃথক কার্ড কিংবা অ্যাডাপ্টার হতে পারে।
\end{itemize}
\end{keypointbox}
\begin{examplebox}{বাস্তব উদাহরণ}
ডেস্কটপ কম্পিউটারের Ethernet পোর্ট বা ল্যাপটপের Wi-Fi অ্যাডাপ্টার NIC-এর উদাহরণ।
\end{examplebox}}{\chTwoImage{chapters/Chapter02/images/nic_connection_model.png}{NIC দিয়ে কম্পিউটারকে নেটওয়ার্কে যুক্ত করা}}


\subsection{Modem}

\begin{definitionbox}{Modem}
ডেটা অনুযায়ী ক্যারিয়ার সংকেতের বৈশিষ্ট্য পরিবর্তন বা Modulation এবং প্রাপ্ত ক্যারিয়ার থেকে ডেটা উদ্ধার বা Demodulation করে যে যন্ত্র, তাকে মডেম বলে। পুরোনো টেলিফোন মডেমের কাজ ডিজিটাল–অ্যানালগ রূপান্তর দিয়ে বোঝানো হয়। আধুনিক সব ইন্টারফেসে একই রূপান্তর ঘটে না।
\end{definitionbox}

Modulator ও Demodulator শব্দ দুটি থেকে Modem নাম এসেছে। মডেম ISP-এর সংযোগ থেকে সংকেত গ্রহণ করে ব্যবহারকারীর ডিভাইস বা রাউটারের উপযোগী ডেটা সংকেত তৈরি করে।

\chTwoPair{\begin{keypointbox}{Modem-এর কাজ}
\begin{itemize}
\item সংকেতের Modulation ও Demodulation করে।
\item ISP-এর সংযোগ ও স্থানীয় নেটওয়ার্ক বা ডিভাইসের মধ্যে যোগাযোগ তৈরি করে।
\item ব্রডব্যান্ড, DSL, কেবল বা ফাইবার সংযোগে প্রযুক্তি অনুযায়ী মডেম বা ONT ব্যবহার করা হয়।
\item ইন্টারনেট সংযোগ চালু করতে গুরুত্বপূর্ণ।
\end{itemize}
\end{keypointbox}
\begin{examplebox}{বাস্তব উদাহরণ}
বাসার ব্রডব্যান্ড সংযোগে ISP-এর সংকেত গ্রহণ করে রাউটার বা কম্পিউটারে ইন্টারনেট ব্যবহারের সুযোগ দিতে সংযোগের ধরন অনুযায়ী মডেম বা ONT ব্যবহার করা হয়।
\end{examplebox}}{\chTwoImage{chapters/Chapter02/images/modem_internet_connection_model.png}{মডেম দিয়ে ইন্টারনেট সংযোগ}}


\subsection{Hub}

\begin{definitionbox}{Hub}
একাধিক ডিভাইসকে LAN-এ সংযুক্ত করে এক পোর্টে প্রাপ্ত ডেটাবাহী সংকেত আগত পোর্ট ছাড়া অন্য পোর্টগুলোতে পাঠায় যে যন্ত্র, তাকে হাব বলে।
\end{definitionbox}

হাব সাধারণত ভৌত স্তরে কাজ করে এবং ঠিকানা দেখে সংকেত বাছাই করে না। সংকেত কোন ডিভাইসের জন্য পাঠানো হয়েছে, হাব তা নির্ধারণ করে না। এক পোর্টে সংকেত পেলে আগত পোর্ট ছাড়া অন্য পোর্টগুলোতে পাঠায়। ফলে অপ্রয়োজনীয় ট্রাফিক ও সংঘর্ষের সম্ভাবনা বাড়ে।

\chTwoPair{\begin{keypointbox}{Hub-এর বৈশিষ্ট্য}
\begin{itemize}
\item একাধিক ডিভাইস সংযুক্ত করতে ব্যবহৃত হয়।
\item আগত পোর্ট ছাড়া অন্য সংযুক্ত পোর্টে সংকেত পাঠায়।
\item MAC Address-এর সারণি ব্যবহার করে না।
\item ঠিকানাভিত্তিক বাছাই করে না; অপ্রয়োজনীয় সংকেত ও একই মাধ্যমে সংঘর্ষের সমস্যা থাকে।
\item আধুনিক নেটওয়ার্কে হাবের পরিবর্তে সুইচ বেশি ব্যবহৃত হয়।
\end{itemize}
\end{keypointbox}
\begin{examplebox}{সহজ উদাহরণ}
পুরোনো একটি ল্যাবের ৫টি কম্পিউটার হাবে যুক্ত আছে। একটি কম্পিউটার ডেটা পাঠালে হাব অন্য সব কম্পিউটারের দিকে সেই সংকেত পাঠাবে।
\end{examplebox}}{\chTwoImage{chapters/Chapter02/images/hub_broadcast_model.png}{হাব আগত পোর্ট ছাড়া অন্য পোর্টে সংকেত পাঠায়}}


\subsection{Switch}

\begin{definitionbox}{Switch}
LAN-এর একাধিক ডিভাইস সংযুক্ত করে MAC Address-এর ভিত্তিতে ডেটা ফ্রেম গন্তব্যে পাঠায় যে নেটওয়ার্ক ডিভাইস, তাকে সুইচ বলে।
\end{definitionbox}


\chTwoDiagramRow{switch}{MAC Table তৈরি ও ব্যবহার}{A থেকে ফ্রেম এলে সুইচ প্রেরকের Source MAC এবং আগত পোর্টের তথ্য শেখে। এরপর Destination MAC খোঁজে। B-এর পোর্ট জানা থাকলে শুধু সেই পোর্টে ফ্রেম পাঠায়। গন্তব্য অজানা হলে বা ফ্রেমটি Broadcast হলে সংশ্লিষ্ট LAN/VLAN-এর অন্য ফরওয়ার্ডিং পোর্টে ছড়িয়ে দেয়। আগত পোর্টে ফেরত পাঠায় না।}
\textbf{সংঘর্ষ ও LAN-এর সীমা:} প্রচলিত Ethernet হাবে সংযুক্ত ডিভাইসগুলো একই Collision Domain ভাগ করে। সুইচের পৃথক ফুল-ডুপ্লেক্স লিংকে সংঘর্ষ ঘটে না। VLAN-সমর্থিত সুইচে একাধিক পৃথক যৌক্তিক LAN থাকতে পারে। একটি সুইচে একটি LAN থাকার ধারণা সাধারণ অবিভক্ত উদাহরণে প্রযোজ্য। Broadcast সংশ্লিষ্ট VLAN-এর মধ্যে সীমাবদ্ধ থাকে।
\textbf{ব্যবহারের উদাহরণ:} ল্যাবের কম্পিউটার থেকে প্রিন্টারে ফ্রেম পাঠাতে সুইচের শেখা পোর্টের সারণি কাজে লাগে। সুইচ প্রেরক কম্পিউটারের MAC তৈরি করে না। Source MAC ফ্রেমে আগে থেকেই থাকে।
\subsection{Router}

\begin{definitionbox}{Router}
দুই বা ততোধিক নেটওয়ার্ক সংযুক্ত করে IP Address-এর ভিত্তিতে ডেটা প্যাকেট পাঠানোর উপযুক্ত পথ নির্ধারণ করে যে ডিভাইস, তাকে রাউটার বলে।
\end{definitionbox}

রাউটার LAN-কে WAN বা ইন্টারনেটের সঙ্গে যুক্ত করে। রাউটিং টেবিলের তথ্য অনুযায়ী প্যাকেট পাঠানোর পরবর্তী পথ নির্ধারণ করে। বাসার Wi-Fi রাউটার সাধারণত স্থানীয় ডিভাইসগুলোকে ইন্টারনেট ব্যবহারের সুযোগ দেয়।

\chTwoPair{\begin{keypointbox}{Router-এর বৈশিষ্ট্য}
\begin{itemize}
\item একাধিক নেটওয়ার্ক সংযুক্ত করে।
\item IP Address-এর ভিত্তিতে প্যাকেট পাঠায়।
\item রাউটিং টেবিলের তথ্য ও মানদণ্ড অনুযায়ী উপযুক্ত পথ নির্ধারণ করে।
\item LAN-কে ইন্টারনেটে যুক্ত করে।
\item অনেক রাউটারে NAT, DHCP ও ফায়ারওয়ালের সুবিধা থাকে।
\end{itemize}
\end{keypointbox}
\begin{examplebox}{বাস্তব উদাহরণ}
বাসার Wi-Fi রাউটার ল্যাপটপ, মোবাইল ফোন ও স্মার্ট টিভিকে স্থানীয় নেটওয়ার্কে যুক্ত করে। মডেম বা ISP-এর সংযোগ দিয়ে ইন্টারনেট ব্যবহারের সুযোগ দেয়।
\end{examplebox}}{\chTwoImage{chapters/Chapter02/images/router_connects_networks_model.png}{রাউটার ভিন্ন নেটওয়ার্ক সংযুক্ত করে}}


\subsection{Repeater}

\begin{definitionbox}{Repeater}
দুর্বল বা বিকৃত সংকেত গ্রহণ করে পুনর্গঠন ও শক্তিশালী করে পরবর্তী অংশে পাঠায় যে নেটওয়ার্ক ডিভাইস, তাকে রিপিটার বলে। সংকেত পুনর্গঠন সাধারণ সংকেত বর্ধনের সমার্থক নয়।
\end{definitionbox}

দীর্ঘ দূরত্ব অতিক্রম করতে গিয়ে ক্ষয়ের কারণে সংকেত দুর্বল হয়। রিপিটার সংকেত নতুন করে তৈরি ও শক্তিশালী করে নেটওয়ার্কের কার্যকর দূরত্ব বাড়াতে সাহায্য করে। তারযুক্ত ও তারবিহীন—উভয় নেটওয়ার্কে এই ধারণা ব্যবহৃত হয়।

\textbf{পুনর্গঠন ও বর্ধনের পার্থক্য:} ডিজিটাল রিপিটার দুর্বল পালস থেকে বিট শনাক্ত করে নতুন পালস তৈরি করে। এতে সংকেতের সময়গত মিল ও আকৃতি পুনর্গঠনে সাহায্য হয়। সাধারণ অ্যাম্প্লিফায়ার সংকেতের সঙ্গে নয়েজও বাড়াতে পারে। হারিয়ে যাওয়া বা ভুল শনাক্ত করা বিট রিপিটার নিশ্চিতভাবে উদ্ধার করতে পারে না।\par
\chTwoPair{\begin{keypointbox}{Repeater-এর কাজ}
\begin{itemize}
\item দুর্বল সংকেত পুনর্গঠন করে।
\item নেটওয়ার্কের কার্যকর দূরত্ব বা সেবার আওতা বাড়ায়।
\item ভৌত স্তরে কাজ করে।
\item ডেটা বাছাই বা রাউটিং করে না।
\item দীর্ঘ কেবল সংযোগ বা Wi-Fi-এর আওতা বাড়াতে ব্যবহার করা যায়।
\end{itemize}
\end{keypointbox}
\begin{examplebox}{বাস্তব উদাহরণ}
বড় ভবনের কিছু স্থানে Wi-Fi সংকেত দুর্বল হলে রিপিটার বা রেঞ্জ এক্সটেন্ডার দিয়ে সেবার আওতা বাড়ানো যায়।
\end{examplebox}}{\chTwoImage{chapters/Chapter02/images/repeater_signal_regeneration.png}{রিপিটার দুর্বল ডিজিটাল সংকেত পুনর্গঠন করে}}


\subsection{Bridge}

\begin{definitionbox}{Bridge}
একই ধরনের নেটওয়ার্কের দুই বা ততোধিক অংশ সংযুক্ত করে MAC Address-এর ভিত্তিতে ট্রাফিক বাছাই করে যে ডিভাইস, তাকে ব্রিজ বলে।
\end{definitionbox}

ব্রিজ LAN-এর অংশগুলো সংযুক্ত করে। MAC Address দেখে কোন অংশে ফ্রেম পাঠাতে হবে, তা নির্ধারণ করে। এতে অপ্রয়োজনীয় ট্রাফিক কমে এবং নেটওয়ার্কের কার্যকারিতা বাড়তে পারে। আধুনিক সুইচকে একাধিক পোর্টবিশিষ্ট ব্রিজ হিসেবেও বোঝানো হয়।

\chTwoPair{\begin{keypointbox}{Bridge-এর বৈশিষ্ট্য}
\begin{itemize}
\item LAN-এর অংশগুলো সংযুক্ত করে।
\item MAC Address দেখে ফ্রেম বাছাই করে।
\item অপ্রয়োজনীয় ট্রাফিক কমায়।
\item ডেটা লিংক স্তরে কাজ করে।
\item ছোট নেটওয়ার্কের অংশ পরিচালনায় ব্যবহার করা যায়।
\end{itemize}
\end{keypointbox}
\begin{examplebox}{সহজ উদাহরণ}
অফিসের হিসাব ও প্রশাসন শাখা একই LAN-এর দুটি অংশ হলে ব্রিজ দিয়ে অংশ দুটি সংযুক্ত এবং ট্রাফিক নিয়ন্ত্রণ করা যায়।
\end{examplebox}}{\chTwoImage{chapters/Chapter02/images/bridge_network_segmentation.png}{ব্রিজ LAN-এর অংশগুলো সংযুক্ত করে}}


\subsection{Gateway}

\begin{definitionbox}{Gateway}
ভিন্ন প্রোটোকল, আর্কিটেকচার বা ধরনের নেটওয়ার্কের মধ্যে ডেটা যোগাযোগ সম্ভব করে যে ডিভাইস বা নোড, তাকে গেটওয়ে বলে।
\end{definitionbox}

প্রোটোকল রূপান্তরের কাজ করা গেটওয়েকে Protocol Converter বলা হয়। রাউটার সাধারণত IP নেটওয়ার্কের মধ্যে প্যাকেটের পথ নির্ধারণ করে। গেটওয়ে ভিন্ন ব্যবস্থা বা প্রোটোকলের মধ্যে প্রয়োজনীয় রূপান্তর করতে পারে। বাসার নেটওয়ার্কে রাউটার সাধারণত Default Gateway-এর ভূমিকা নেয়।

\chTwoPair{\begin{keypointbox}{Gateway-এর বৈশিষ্ট্য}
\begin{itemize}
\item ভিন্ন নেটওয়ার্ক বা প্রোটোকলের ব্যবস্থা সংযুক্ত করে।
\item প্রোটোকল রূপান্তর করতে পারে।
\item নেটওয়ার্কের প্রবেশ ও নির্গমনপথের ভূমিকা নিতে পারে।
\item প্রোটোকল রূপান্তরের কাজ রাউটিংয়ের তুলনায় জটিল হতে পারে।
\item ইন্টারনেট প্রবেশ বা প্রতিষ্ঠানের বিভিন্ন নেটওয়ার্ক সমন্বয়ে গুরুত্বপূর্ণ।
\end{itemize}
\end{keypointbox}
\begin{examplebox}{বাস্তব উদাহরণ}
স্থানীয় নেটওয়ার্ক থেকে ইন্টারনেটে ডেটা পাঠাতে Default Gateway ব্যবহৃত হয়। আবার ভিন্ন প্রোটোকলভিত্তিক দুটি ব্যবস্থা যুক্ত করতে প্রোটোকল রূপান্তরকারী গেটওয়ে প্রয়োজন হতে পারে।
\end{examplebox}}{\chTwoImage{chapters/Chapter02/images/gateway_protocol_conversion.png}{গেটওয়ে ভিন্ন নেটওয়ার্ক বা প্রোটোকল ব্যবস্থা যুক্ত করে}}


\subsection{Access Point}

\begin{definitionbox}{Access Point}
তারবিহীন ডিভাইসকে তারযুক্ত স্থানীয় নেটওয়ার্ক বা LAN-এ যুক্ত হওয়ার সুযোগ দেয় যে নেটওয়ার্ক ডিভাইস, তাকে Access Point বা AP বলে।
\end{definitionbox}

অ্যাকসেস পয়েন্ট Wi-Fi সংকেত দিয়ে মোবাইল ফোন, ল্যাপটপ ও ট্যাবলেটকে নেটওয়ার্ক ব্যবহারের সুযোগ দেয়। বাসার অনেক রাউটারে অ্যাকসেস পয়েন্ট যুক্ত থাকে।

\chTwoPair{\begin{keypointbox}{Access Point-এর কাজ}
\begin{itemize}
\item তারবিহীন LAN তৈরি করে।
\item Wi-Fi ডিভাইসকে নেটওয়ার্ক ব্যবহারের সুযোগ দেয়।
\item তারযুক্ত নেটওয়ার্ক ও তারবিহীন ব্যবহারকারীর সংযোগস্থল হিসেবে কাজ করে।
\item বিদ্যালয়, অফিস, গ্রন্থাগার, ক্যাম্পাস ও বাসার নেটওয়ার্কে ব্যবহৃত হয়।
\item সেবার আওতা বাড়াতে একাধিক অ্যাকসেস পয়েন্ট ব্যবহার করা যায়।
\end{itemize}
\end{keypointbox}
\begin{examplebox}{বাস্তব উদাহরণ}
কলেজ ক্যাম্পাসের বিভিন্ন স্থানে অ্যাকসেস পয়েন্ট বসিয়ে শিক্ষার্থীদের Wi-Fi সংযোগ দেওয়া যায়। নেটওয়ার্কে ইন্টারনেট সংযোগ থাকলে তাঁরা ইন্টারনেটও ব্যবহার করতে পারেন।
\end{examplebox}}{\chTwoImage{chapters/Chapter02/images/access_point_wireless_lan.png}{অ্যাকসেস পয়েন্ট বেতার যন্ত্রকে নেটওয়ার্কে যুক্ত করে}}


\subsection{Firewall}

\begin{definitionbox}{Firewall}
পূর্বনির্ধারিত নিরাপত্তানীতি অনুযায়ী আগত ও বহির্গামী নেটওয়ার্ক ট্রাফিক পর্যবেক্ষণ ও বাছাই করে যে নিরাপত্তা ডিভাইস বা সফটওয়্যার ব্যবস্থা, তাকে ফায়ারওয়াল বলে।
\end{definitionbox}

ফায়ারওয়াল অননুমোদিত প্রবেশ, ক্ষতিকর ট্রাফিক, সন্দেহজনক সংযোগ ও অবাঞ্ছিত ডেটা আটকাতে সাহায্য করে। এটি হার্ডওয়্যার, সফটওয়্যার বা রাউটারের অন্তর্ভুক্ত ব্যবস্থা হতে পারে।

\chTwoPair{\begin{keypointbox}{Firewall-এর কাজ}
\begin{itemize}
\item নেটওয়ার্ক ট্রাফিক পর্যবেক্ষণ করে।
\item নিরাপত্তানীতি অনুযায়ী ট্রাফিক অনুমোদন বা বাধা দেয়।
\item অননুমোদিত প্রবেশ কমাতে সাহায্য করে।
\item ম্যালওয়্যার, আক্রমণ বা অবাঞ্ছিত সংযোগের ঝুঁকি কমাতে সাহায্য করে।
\item বাসার রাউটার, অপারেটিং সিস্টেম ও প্রতিষ্ঠানের নেটওয়ার্কে ফায়ারওয়াল ব্যবহৃত হয়।
\end{itemize}
\end{keypointbox}
\begin{examplebox}{বাস্তব উদাহরণ}
অফিসের নেটওয়ার্কে ফায়ারওয়াল ব্যবহার করে বাইরের অননুমোদিত ব্যবহারকারী বা ক্ষতিকর ট্রাফিকের সার্ভারে প্রবেশ নিয়ন্ত্রণ করা যায়।
\end{examplebox}}{\chTwoImage{chapters/Chapter02/images/firewall_filtering_model.png}{ফায়ারওয়াল ট্রাফিক বাছাই করে}}


\subsection{Home Network-এ Network Device-এর ব্যবহার}

বাসার সাধারণ নেটওয়ার্কে মডেম বা ONT, রাউটার, সুইচ, অ্যাকসেস পয়েন্ট ও NIC একসঙ্গে কাজ করতে পারে।


\chTwoPair{\begin{keypointbox}{Home network flow}
\begin{itemize}
\item ISP-এর সংকেত মডেম বা ONT-এ আসে।
\item রাউটার স্থানীয় ডিভাইসগুলোকে ইন্টারনেট ব্যবহারের সুযোগ দেয়।
\item রাউটারের অন্তর্ভুক্ত অ্যাকসেস পয়েন্ট Wi-Fi সংকেত দেয়।
\item ডেস্কটপ কম্পিউটার NIC বা Ethernet পোর্ট দিয়ে রাউটারে যুক্ত হয়।
\item মোবাইল ফোন ও ল্যাপটপ তারবিহীন NIC দিয়ে Wi-Fi নেটওয়ার্কে যুক্ত হয়।
\end{itemize}
\end{keypointbox}}{\chTwoImage{chapters/Chapter02/images/home_network_devices_layout.png}{বাসার নেটওয়ার্কে ডিভাইসের ব্যবহার}}


\subsection{School Computer Lab Network}

বিদ্যালয় বা কলেজের কম্পিউটার ল্যাবে সাধারণত সুইচ, রাউটার, মডেম বা ONT, সার্ভার, প্রিন্টার ও কম্পিউটার ব্যবহৃত হয়।


\chTwoImage{chapters/Chapter02/images/school_lab_network_devices_layout.png}{বিদ্যালয়ের কম্পিউটার ল্যাবের নেটওয়ার্ক বিন্যাস}


\subsection{ডিভাইসের পার্থক্য: কোন সূত্রে চিনবে?}
\begin{longtable}{|p{\dimexpr .24\linewidth-7pt\relax}|p{\dimexpr .76\linewidth-7pt\relax}|}\hline
\textbf{তুলনা} & \textbf{সিদ্ধান্তের মূল পার্থক্য}\\\hline
\endfirsthead\hline\textbf{তুলনা} & \textbf{সিদ্ধান্তের মূল পার্থক্য}\\\hline\endhead
Hub / Switch & হাব ঠিকানা শেখে না। সুইচ MAC Table দিয়ে পরিচিত Unicast-এর অপ্রয়োজনীয় ট্রাফিক কমায়। গন্তব্য অজানা হলে সুইচও ফ্রেম ছড়িয়ে দিতে পারে।\\\hline
Bridge / Switch & দুটিই ব্রিজিং ও MAC-ভিত্তিক বাছাই করে। সুইচ মূলত একাধিক পোর্টবিশিষ্ট ব্রিজ। পুরোনো ব্রিজে সাধারণত পোর্ট ও বহনক্ষমতা কম ছিল। তবে কোনো শ্রেণির গতি সর্বজনীন স্থির মান নয়।\\\hline
Repeater / Bridge & রিপিটার সংকেত পুনর্গঠন করে। ব্রিজ MAC দেখে ফ্রেম পাঠানোর প্রয়োজনীয় অংশ বেছে নেয়। সংকেতের দূরত্ব বাড়ানো এবং ট্রাফিক বাছাই আলাদা কাজ।\\\hline
Switch / Router & সাধারণ Layer-2 সুইচ একই LAN/VLAN-এর ফ্রেম নিয়ে কাজ করে। রাউটার ভিন্ন IP নেটওয়ার্কের প্যাকেটের পথ নির্ধারণ করে।\\\hline
Router / Protocol Gateway & রাউটার IP-ভিত্তিক পথ নির্ধারণ করে; প্রোটোকল গেটওয়ে প্রোটোকল বা অ্যাপ্লিকেশনের রূপান্তর করে। রাউটার Default Gateway হতে পারে। সেই ভূমিকা মানেই প্রোটোকল রূপান্তর নয়।\\\hline
Modem / Router & মডেমের মূল কাজ Modulation/Demodulation; রাউটারের কাজ IP নেটওয়ার্কের মধ্যে পথ নির্ধারণ। একই যন্ত্রে দুটি সুবিধা থাকতে পারে।\\\hline
NIC / Modem & NIC হোস্টের নেটওয়ার্ক ইন্টারফেস। মডেম ক্যারিয়ার ও ডেটার রূপান্তরের যন্ত্র। সব NIC মডেম নয়।\\\hline
\end{longtable}
\chTwoImage{chapters/Chapter02/images/hub_switch_router_comparison.png}{সংকেত, MAC ও IP-ভিত্তিক কাজের পার্থক্য}
\subsection{কোন পরিস্থিতিতে কোন device ব্যবহার হবে?}

\begingroup
\begin{longtable}{|p{\dimexpr 0.4070\linewidth-2\tabcolsep-\arrayrulewidth\relax}|p{\dimexpr 0.2714\linewidth-2\tabcolsep-\arrayrulewidth\relax}|p{\dimexpr 0.3166\linewidth-2\tabcolsep-\arrayrulewidth\relax}|}
\hline
\textbf{পরিস্থিতি} & \textbf{উপযুক্ত device} & \textbf{কারণ} \\
\hline
\endfirsthead
\hline
\textbf{পরিস্থিতি} & \textbf{উপযুক্ত device} & \textbf{কারণ} \\
\hline
\endhead

কম্পিউটারকে নেটওয়ার্কে যুক্ত করা & NIC & কম্পিউটারের নেটওয়ার্ক ইন্টারফেস প্রয়োজন \\
\hline
ISP-এর সংযোগে ইন্টারনেট নেওয়া & Modem/ONT & প্রযুক্তি অনুযায়ী ISP-এর সংকেত স্থানীয় যন্ত্র বা নেটওয়ার্কের উপযোগী করে \\
\hline
ল্যাবের অনেক কম্পিউটার যুক্ত করা & Switch & LAN-এর ডিভাইসগুলো দক্ষতার সঙ্গে সংযুক্ত করে \\
\hline
দুর্বল সংকেত দূরে পাঠানো & Repeater & সংকেত পুনর্গঠন করে কার্যকর দূরত্ব বাড়ায় \\
\hline
LAN-এর দুটি অংশ যুক্ত করা & Bridge & অংশগুলো সংযুক্ত করে ফ্রেম বাছাই করে \\
\hline
LAN-কে ইন্টারনেটে যুক্ত করা & Router & ভিন্ন নেটওয়ার্ক যুক্ত করে পথ নির্ধারণ করে \\
\hline
ভিন্ন প্রোটোকলের ব্যবস্থা যুক্ত করা & Gateway & প্রোটোকল রূপান্তর করতে পারে \\
\hline
তারবিহীন ইন্টারনেট ব্যবহারের সুযোগ দেওয়া & Access Point & Wi-Fi ডিভাইস যুক্ত করে; নেটওয়ার্কে ইন্টারনেট সংযোগ প্রয়োজন \\
\hline
অননুমোদিত প্রবেশ নিয়ন্ত্রণ & Firewall & নিরাপত্তানীতি অনুযায়ী ট্রাফিক বাছাই করে \\
\hline
\end{longtable}
\endgroup

\par\Needspace{6\baselineskip}
\subsection{Hub ও Bridge-এর ঐতিহাসিক ধরন}
\textbf{পুরোনো হাবের ধরন:} Passive Hub মূলত সংযোগ ভাগ করে; সংকেত পুনর্গঠন করে না। Active Hub বিদ্যুৎ ব্যবহার করে সংকেত পুনর্গঠন করে। Intelligent Hub-এ ব্যবস্থাপনার সুবিধা থাকতে পারে। আধুনিক ল্যাবে সাধারণত সুইচ ব্যবহার করা হয়।
\textbf{ব্রিজের ব্যবহারভেদ:} Local Bridge কাছের অংশ এবং Remote Bridge দূরের অংশ যুক্ত করে। Wireless Bridge বেতার সংযোগ দিয়ে নেটওয়ার্কের অংশ সংযুক্ত করে। এই নামগুলো মাধ্যম ও অবস্থানের পরিচয় দেয়। MAC-এর ভিত্তিতে ব্রিজের মূল কাজ বদলায় না।

\subsection{Modem-এর অবস্থান ও Router-এর সহায়ক কাজ}
\begin{multicols}{2}
\textbf{Internal ও External Modem:} Internal Modem কম্পিউটারের ভেতরে কার্ড বা মডিউল হিসেবে থাকে। External Modem বাইরে থাকে এবং কেবল বা পোর্ট দিয়ে যুক্ত হয়। পুরোনো ডেস্কটপের ডায়াল-আপ মডেম কার্ড প্রথমটির এবং বাইরে রাখা USB মডেম দ্বিতীয়টির উদাহরণ। অবস্থান ও সংযোগের ভিত্তিতে এই শ্রেণি করা হয়। Modulation ও Demodulation-এর মূল কাজ একই থাকে।\par
\textbf{একটি Public IP-তে একাধিক সংযোগ:} NAT (Network Address Translation) IP Address রূপান্তর করে। PAT (Port Address Translation) বা NAT Overload-এ IP-এর পাশাপাশি পৃথক TCP/UDP Port Number দিয়ে এক Public IP-তে একাধিক কম্পিউটারের সংযোগ আলাদা রাখা হয়। ল্যাবের দুটি কম্পিউটার একই ওয়েবসাইট খুললেও রূপান্তরের সারণি অনুযায়ী উত্তর সঠিক কম্পিউটারে ফিরে যায়।\par
\end{multicols}
\textbf{DHCP (Dynamic Host Configuration Protocol):} DHCP সার্ভার স্বয়ংক্রিয়ভাবে IP Address ও প্রয়োজনীয় নেটওয়ার্ক কনফিগারেশন দেয়। বাসার রাউটারে এই সুবিধা থাকলে নতুন ফোনে ঠিকানা হাতে লিখতে হয় না। Routing, NAT/PAT ও DHCP আলাদা কাজ। সব রাউটারে সব সুবিধা থাকা আবশ্যক নয়। Default Gateway হিসেবে কাজ করলেই যন্ত্রটি প্রোটোকল রূপান্তরকারী হয়ে যায় না।
\textbf{DSL ও ADSL:} DSL = Digital Subscriber Line। টেলিফোনের তামার জোড়া তারে উপযুক্ত মডেম ও সংযোগযন্ত্র দিয়ে ডেটা সেবা দেওয়া যায়। কথার ও ডেটার সংকেত পৃথক ফ্রিকোয়েন্সি অঞ্চলে থাকলে ফিল্টার বা স্প্লিটার দিয়ে একই সংযোগে দুটিই ব্যবহার করা সম্ভব। ADSL (Asymmetric DSL)-এ ডাউনস্ট্রিম ও আপস্ট্রিমের বহনক্ষমতা অসমান; সাধারণত ডাউনলোডের ক্ষমতা বেশি। এটি সংযোগের প্রযুক্তি। Internal/External মডেমের অবস্থান এবং Star/Bus সংযোগের বিন্যাস বোঝায়।

\subsection{Modem, ONT ও Gateway নিয়ে বিভ্রান্তি}
\begin{tabularx}{\linewidth}{|p{.21\linewidth}|X|X|}
\hline
\textbf{নাম} & \textbf{মূল কাজ} & \textbf{একটির সঙ্গে অন্যটি গুলিয়ে ফেলবে না}\\\hline
Modem & ডেটা অনুযায়ী বাহক সংকেত বদলানো এবং প্রাপ্ত সংকেত থেকে ডেটা উদ্ধার করা; অর্থাৎ Modulation/Demodulation। & পুরোনো টেলিফোন মডেম দিয়ে ডিজিটাল–অ্যানালগ রূপান্তর বোঝানো যায়। আধুনিক সব ইন্টারনেট ইন্টারফেস একইভাবে রূপান্তর করে না।\\\hline
ONT & Optical Network Terminal: ফাইবার সংযোগের প্রান্তে আলোকসংকেত গ্রহণ করে স্থানীয় ডেটা ইন্টারফেস দেয়। & FTTH-তে ব্যবহৃত হয়। টেলিফোনের অ্যানালগ মডেম ও ONT একই যন্ত্র নয়।\\\hline
Default Gateway & স্থানীয় নেটওয়ার্কের বাইরের গন্তব্যে পাঠানোর পরবর্তী যন্ত্র; সাধারণত রাউটার। & কম্পিউটারের গেটওয়ে সেটিং মানেই প্রোটোকল রূপান্তরকারী নয়।\\\hline
Protocol Gateway & ভিন্ন প্রোটোকল বা অ্যাপ্লিকেশন ব্যবস্থার মধ্যে প্রয়োজনীয় রূপান্তর করে। & প্রোটোকল রূপান্তরের সাহায্যে ভিন্ন ব্যবস্থার মধ্যে যোগাযোগ সম্ভব হয়।\\\hline
\end{tabularx}
\textbf{একটি যন্ত্রে একাধিক কাজ:} বাসার তারবিহীন রাউটারে Routing, Ethernet Switch ও Wi-Fi Access Point-এর সুবিধা একত্রে থাকতে পারে। ফায়ারওয়ালও থাকতে পারে। তাই প্রতিটি কাজের জন্য পৃথক যন্ত্র আঁকা সবসময় প্রয়োজন নয়। তবে রাউটার নাম থেকেই সব মডেলে মডেম বা ONT আছে বলে ধরে নেওয়া যায় না।
\textbf{নিরাপত্তা:} ফায়ারওয়াল নির্ধারিত নীতি অনুযায়ী ট্রাফিক অনুমোদন বা বাধা দেয়। শক্তিশালী পাসওয়ার্ড, হালনাগাদ, ব্যাকআপ ও সচেতন ব্যবহারের প্রয়োজনও থাকে। সুইচ অপ্রয়োজনীয় ফ্রেম চলাচল কমায়। তবে শুধু সুইচ ব্যবহারের কারণে গোপন তথ্য স্বয়ংক্রিয়ভাবে এনক্রিপ্ট হয় না।


\textbf{যন্ত্রের পোর্ট ও দূরত্বের সংখ্যা:} ৮/১৬/২৪/৪৮-Port সুইচ পরিচিত কনফিগারেশন; এগুলো পোর্টসংখ্যার সংজ্ঞাগত সর্বোচ্চ সীমা নয়। ল্যাবের ১২ PC ও ১ Uplink-এর জন্য অন্তত ১৩টি উপযুক্ত পোর্ট প্রয়োজন। ডেটা হার ও অন্যান্য সুবিধা মিললে ১৬-Port মডেল নেওয়া যায়। রিপিটার বসানোর ব্যবধান কেবল বা বেতার মানের ওপর নির্ভর করে। ১৮৫ m পুরোনো 10BASE2 অংশের দৈর্ঘ্য; সব রিপিটারের জন্য একই ব্যবধান নয় (অংশ \ref{ch2:wired-numeric})।

\textbf{রাউটার ও গেটওয়ের ইতিহাস:} ১৯৭০-এর দশকে পৃথক নেটওয়ার্ক সংযুক্ত করার পরীক্ষায় Virginia Strazisar-এর প্রাথমিক TCP/IP Gateway/Router সফটওয়্যার গুরুত্বপূর্ণ অবদান রাখে। হাব, ব্রিজ, রিপিটার ও গেটওয়ে বিভিন্ন নকশায় বিকশিত হয়েছে। সব শ্রেণির জন্য একক উদ্ভাবক ও সাল নির্ধারণ নির্ভরযোগ্য নয়। Ethernet-এর উৎপত্তি নেটওয়ার্কের ইতিহাসে আছে। আধুনিক Ethernet সুইচ ও পুরোনো ভাগাভাগি করা বাস একই যন্ত্র নয়।

% Original application; device mechanisms remain in their theory sections.
\subsection{বাস্তব সমস্যা: Lab চলে, Internet চলে না}
\chTwoPair{\textbf{সমস্যা:} ল্যাবের কম্পিউটারগুলো সুইচ দিয়ে স্থানীয় ফাইল ব্যবহার করতে পারছে। কিন্তু ইন্টারনেট সংযোগ পাওয়া যাচ্ছে না।\par
\textbf{সমাধান:} কম্পিউটারের IP কনফিগারেশন, Default Gateway, রাউটার ও ঊর্ধ্বমুখী সংযোগ পরীক্ষা করতে হবে। স্থানীয় সংযোগ সচল থাকায় শুধু সুইচ বদলালে সমাধান নাও হতে পারে।}{\textbf{নতুন PC-র ক্ষেত্রে:} সুইচের MAC Table-এ প্রাপকের ঠিকানা শুরুতে অজানা থাকতে পারে। সুইচের অংশে বর্ণিত ঠিকানা শেখা ও ফরওয়ার্ড করার নিয়ম এখানেও প্রযোজ্য। সার্ভার বা অ্যাপ্লিকেশন অচল কি না, তা পৃথকভাবে পরীক্ষা করতে হবে।}

\textbf{LAN/WAN-এ ডিভাইসের ভূমিকা:} স্থানীয় ডিভাইস সংযুক্ত করতে হাব বা সুইচ এবং নেটওয়ার্কগুলোর মধ্যে সংযোগ দিতে রাউটার বা গেটওয়ে ব্যবহার করা হয়। NIC-ও LAN-এ ডিভাইস যুক্ত করে। তাই প্রশ্নে শুধু LAN Device উল্লেখ থাকলে অন্য সব যন্ত্র বাদ দেওয়া যায় না। MAC ফ্রেম, IP প্যাকেট, প্রোটোকল রূপান্তর ও সংকেতের কাজ দিয়ে পার্থক্য নির্ধারণ করতে হয়।


\subsection{সাধারণ ভুল ও সংশোধন}

\smallhead{১. Switch ও router একই মনে করা}
সুইচ একই LAN-এর ডিভাইসগুলোর মধ্যে ফ্রেম পাঠায়। রাউটার ভিন্ন নেটওয়ার্ক সংযুক্ত করে এবং IP Address-এর ভিত্তিতে প্যাকেটের পথ নির্ধারণ করে। তাই সুইচ ও রাউটারের কাজ আলাদা।

\smallhead{২. Hub ও switch দেখতে একই বলে কাজও একই মনে করা}
হাব ও সুইচ দেখতে একই রকম হতে পারে। হাব আগত পোর্ট ছাড়া অন্য পোর্টে সংকেত পাঠায়। সুইচ গন্তব্যের MAC Address দেখে পরিচিত গন্তব্যের উপযুক্ত পোর্টে ফ্রেম পাঠায়।

\smallhead{৩. Modem ও router এক জিনিস মনে করা}
মডেম ISP-এর সংকেত রূপান্তর করে। রাউটার স্থানীয় ডিভাইসগুলোকে অন্য নেটওয়ার্ক বা ইন্টারনেটে যোগাযোগের সুযোগ দেয়। দুটি কাজ পৃথক।

\smallhead{৪. Gateway ও router concept গুলিয়ে ফেলা}
গেটওয়ে ও রাউটারের সুবিধা একই যন্ত্রে থাকতে পারে; তবে ধারণা দুটি আলাদা। গেটওয়ে নেটওয়ার্কের প্রবেশ ও নির্গমনপথের ভূমিকা নিতে পারে। প্রোটোকল গেটওয়ে ভিন্ন ব্যবস্থার মধ্যে প্রয়োজনীয় রূপান্তরও করতে পারে।

\subsection{MCQ অনুশীলন}

\par
\textbf{MCQ}

\begin{multicols}{2}\begin{enumerate}
\item\chTwoMCQBlock{কম্পিউটারকে নেটওয়ার্কে যুক্ত করতে কোন উপাদান প্রয়োজন?
\par ক. NIC \quad খ. মনিটর \par গ. স্ক্যানার \quad ঘ. স্পিকার }
\item\chTwoMCQBlock{কোন ডিভাইস MAC Address-এর ভিত্তিতে ফ্রেম ফরওয়ার্ড করে?
\par ক. হাব \quad খ. সুইচ \par গ. মডেম \quad ঘ. কিবোর্ড }
\item\chTwoMCQBlock{কোন ডিভাইস IP Address-এর ভিত্তিতে প্যাকেটের পথ নির্ধারণ করে?
\par ক. রাউটার \quad খ. হাব \par গ. রিপিটার \quad ঘ. মাউস }
\item\chTwoMCQBlock{হাব প্রাপ্ত সংকেত কোথায় পাঠায়?
\par ক. শুধু নির্দিষ্ট পোর্টে \par খ. আগত পোর্ট ছাড়া অন্য সংযুক্ত পোর্টে \par গ. শুধু ইন্টারনেটে \quad ঘ. কোথাও পাঠায় না }
\item\chTwoMCQBlock{দুর্বল ডিজিটাল সংকেত পুনর্গঠনে কোন ডিভাইস ব্যবহার করা হয়?
\par ক. রিপিটার \quad খ. প্রিন্টার \par গ. স্ক্যানার \quad ঘ. মনিটর }
\item\chTwoMCQBlock{ভিন্ন প্রোটোকলের ব্যবস্থা যুক্ত করতে কোন ডিভাইস উপযোগী?
\par ক. গেটওয়ে \quad খ. কিবোর্ড \par গ. স্পিকার \quad ঘ. প্রজেক্টর }
\item\chTwoMCQBlock{Wi-Fi ডিভাইসকে নেটওয়ার্কে যুক্ত করতে কোনটি ব্যবহৃত হয়?
\par ক. অ্যাকসেস পয়েন্ট \quad খ. প্লটার \par গ. বারকোড রিডার \quad ঘ. জয়স্টিক }
\item\chTwoMCQBlock{নিরাপত্তানীতি অনুযায়ী নেটওয়ার্ক ট্রাফিক বাছাই করে কোনটি?
\par ক. ফায়ারওয়াল \quad খ. হাব \par গ. মাউস \quad ঘ. মনিটর }
\end{enumerate}\end{multicols}

\textbf{উত্তরমালা:}
১. ক, ২. খ, ৩. ক, ৪. খ, ৫. ক, ৬. ক, ৭. ক, ৮. ক

\smallskip
\endgroup
